\chapter{Flux}\label{ch:s1:flows}

Lorsque des investisseurs retirent de l'argent d'un fonds commun de placement, le fonds vend ce qu'il détient~; lorsque de nombreux fonds perdent de l'argent ensemble, ils vendent ensemble les mêmes actions, et
les cours baissent pour des raisons sans rapport avec les entreprises. Coval et Stafford ont mis en évidence cette \omterm{def:s1:flows:fit}{pression sur les prix} dans les transactions de fonds communs de placement entre 1980 et
2003, et ont constaté que celui qui négociait contre les vendeurs contraints était rémunéré pour cela. Lou a agrégé les transactions imposées par les flux des fonds en une mesure pour chaque
action et a relevé un impact sur les prix significatif et temporaire, qui s'inverse les années suivantes~; les flux étant prévisibles, la pression l'est aussi. Ce chapitre
simule deux cents fonds sur le marché synthétique, mesure la pression qu'exercent leurs flux et la part qui s'inverse, et teste les deux transactions que
suggère la littérature~: devancer les flux, et acheter ce que les flux ont forcé à vendre. La construction est \texttt{firm.flowpress}.

\section{Pression sur les prix induite par les flux}

\begin{definition}[Transactions induites par les flux, pression sur les prix]\label{def:s1:flows:fit}
Les \emph{transactions induites par les flux}\index{transactions induites par les flux} d'une action sur une période sont l'ensemble des transactions que les flux des fonds lui imposent~: la somme, sur les fonds, de la part de
l'action que chaque fonds détient, multipliée par le flux du fonds rapporté à ses actifs. La \emph{pression sur les prix}\index{pression sur les prix} est la part d'un mouvement de cours causée par
une demande non informée de ce type, censée s'inverser une fois la demande absorbée.
\end{definition}

Pourquoi une demande non informée devrait-elle faire bouger les cours~? Si d'autres investisseurs étaient prêts à prendre la contrepartie à un prix proche de la valeur, elle ne le ferait pas. Gabaix et
Koijen soutiennent que le marché, dans son ensemble, est inélastique~: les institutions opèrent sous des mandats qui laissent peu de marge pour absorber les flux, et estiment que
1~\$ investi dans le marché boursier fait monter sa valeur d'environ 5~\$. Pour une action isolée, les \omterm{def:s1:flows:fit}{transactions induites par les flux}, telles que les mesurent Coval et Stafford et Lou, en sont la
version mesurée de la même idée.

Les fonds du marché synthétique (\texttt{firm.flowpress}) sont simples. Deux cents fonds détiennent chacun soixante actions, tirées au hasard, et remplacent chaque trimestre un quart de leurs
positions~; ensemble, ils détiennent environ 20~\% de chaque action. Chaque trimestre, le flux d'un fonds représente 5~\% de ses actifs par écart type de son rendement
annuel glissant, relatif à celui des autres fonds, plus un bruit de 3~\%~: les flux courent après la performance. Les \omterm{def:s1:flows:fit}{transactions induites par les flux} d'une action ce trimestre-là (écart type de 0.5~\% de
ses actions) font varier son cours de 1.5 fois leur montant, étalé sur le trimestre, et cette poussée s'inverse avec une demi-vie de 126 jours. Comme les flux courent après
des rendements qui intègrent la pression passée, la pression engendre des flux~: le mécanisme que Lou a proposé pour la persistance de la performance des fonds.

\section{Rachats de fonds communs de placement et ventes forcées}

Mesuré en coupe transversale, le rendement logarithmique de chaque trimestre varie de 1.05 par unité de \omterm{def:s1:flows:fit}{transactions induites par les flux} de ce trimestre~; sur l'année suivante, $-0.69$, de sorte que
66~\% du mouvement du trimestre s'inverse en un an (la décroissance programmée seule reprendrait 75~\% de la poussée de fin de trimestre). Sans pression programmée, les
mêmes pentes valent 0.31 et 0.32~: les fonds qui ont bien performé détiennent des actions qui ont bien performé, et les actions du marché synthétique ont des dérives persistantes, si bien que les flux visent
des actions en dérive, qu'elles les poussent ou non. Toute mesure de la pression des flux doit être lue à la lumière de cette variable confondante.

\begin{omfigure}
\begin{tikzpicture}
\begin{axis}[width=11cm,height=5.4cm,xlabel={\small années depuis le début du trimestre},ylabel={\small décile supérieur moins décile inférieur (\%)},
  tick label style={font=\footnotesize},xmin=0,xmax=3,ymin=-4,ymax=4,grid=major,grid style={gray!25},legend pos=north east,
  legend style={font=\scriptsize},table/col sep=comma]
\addplot[omDef,very thick] table[x=year,y=spread]{figdata/strategies-1/09-flows/path.csv};
\addplot[black!50,thick,dashed] table[x=year,y=no_pressure]{figdata/strategies-1/09-flows/path.csv};
\addplot[black!40,dotted] coordinates {(0,0) (3,0)};
\legend{avec pression des flux,mêmes fonds et mêmes flux sans pression}
\end{axis}
\end{tikzpicture}

\omcaption{\omterm{def:s1:earnings:event}{Étude d'événement} des \omterm{def:s1:flows:fit}{transactions induites par les flux} sur le marché synthétique~: rendement logarithmique cumulé des actions appartenant au décile supérieur des \omterm{def:s1:flows:fit}{transactions induites par les flux}
d'un trimestre moins le décile inférieur, depuis le début du trimestre, avec la pression programmée et sans elle. Données~: \texttt{s1\_flows.event\_path}.}
\label{fig:s1:flows:path}
\end{omfigure}

L'\omterm{def:s1:earnings:event}{étude d'événement} (la \cref{fig:s1:flows:path}) montre la forme décrite par Lou. Le décile supérieur des \omterm{def:s1:flows:fit}{transactions induites par les flux} d'un trimestre devance le décile inférieur de 1.74~\% sur le
trimestre et de 1.91~\% à son maximum, trois jours après la fin du trimestre~; l'écart revient à zéro 261 jours de bourse après le début du trimestre, soit à peu près un an, et finit à $-2.15~\%$ après
trois ans. Sans pression, l'écart vaut 0.15~\% à la fin du trimestre et dérive jusqu'à $-1.41~\%$.

\section{Flux des ETF et des fonds indiciels}

Les fonds communs de placement gèrent leurs flux avec une certaine latitude~; les ETF ne le peuvent pas. Les créations et rachats d'un ETF portent sur des paniers de ses titres, et
les arbitragistes échangent le panier chaque fois que le cours de l'ETF s'écarte de sa valeur. Ben-David, Franzoni et Moussawi, qui utilisent les changements d'appartenance aux indices comme
expérience naturelle, constatent que les actions davantage détenues par des ETF présentent une volatilité nettement plus élevée et une autocorrélation plus négative (les chocs de liquidité dans l'ETF
atteignent ses paniers par l'arbitrage) et rapportent une prime de risque pouvant atteindre 56 points de base par mois. Les flux des fonds indiciels sont les plus mécaniques de tous,
et le chapitre~10 exploite le plus prévisible d'entre eux, le rééquilibrage des indices.

\section{Positions trimestrielles comme signal}

\begin{definition}[Positions 13F]\label{def:s1:flows:13f}
Les \emph{positions 13F}\index{positions 13F} sont les positions trimestrielles sur actions que les gestionnaires d'investissement institutionnels américains dont les encours dépassent un seuil doivent déclarer à
la SEC au moyen du Form 13F, publiées après un délai~; elles constituent la source publique des positions des fonds sur lesquelles sont construits les signaux de flux.
\end{definition}

\begin{dated}{2026-09}{Form 13F}\label{dat:s1:flows:13f}
Les gestionnaires d'investissement institutionnels exerçant un pouvoir discrétionnaire d'investissement sur des titres de la section 13(f) d'un montant de 100~millions de dollars ou plus doivent déposer le Form 13F auprès de la SEC
chaque trimestre, dans les 45 jours suivant la fin du trimestre~; la SEC publie la liste officielle des titres relevant de la section 13(f) peu après chaque trimestre. Les fonds communs
de placement déclarent aussi l'intégralité de leurs portefeuilles sur leurs propres formulaires. Un signal fondé sur les positions voit donc des positions âgées d'un mois et demi à quatre mois et demi.
\end{dated}

Les deux backtests du chapitre respectent ce délai.

\begin{center}
\small
\setlength{\tabcolsep}{4pt}%
\begin{tabular}{@{}lcc@{}}
\toprule
fonds synthétiques, années 3 à 10 & avec pression programmée & sans\\
\midrule
flux anticipés~: ratio de Sharpe avant, après coûts & 0.70, 0.44 & 0.38, 0.12\\
flux anticipés~: rendement annuel après coûts~; rotation & 1.4~\%~; 4.0 & ---\\
inversion après vente forcée~: ratio de Sharpe avant, après coûts & $-0.03$, $-0.35$ & $-0.73$, $-1.06$\\
\bottomrule
\end{tabular}
\end{center}

Le backtest \emph{flux anticipés} prévoit à chaque début de trimestre les \omterm{def:s1:flows:fit}{transactions induites par les flux} du trimestre suivant~: flux issus des rendements glissants des fonds (avec une pente
estimée sur les seuls trimestres passés) multipliés par les positions de fin du trimestre précédent, déjà déclarées. Il achète le décile supérieur et vend le décile inférieur pour le trimestre,
et affiche un ratio de Sharpe net de 0.44. Le backtest \emph{inversion après vente forcée} attend que les transactions d'un trimestre soient connues grâce à la déclaration suivante, 45 jours après la fin du trimestre,
achète les actions que les flux ont le plus vendues et vend celles qu'ils ont achetées, et tient la position un an. Il perd de l'argent~: au moment où les transactions sont connues, une bonne part de
l'inversion a déjà eu lieu, et ce qui reste est contrarié par les dérives que les flux poursuivaient. La version sans pression perd davantage, ce qui représente la
contribution de la pression. L'inversion est réelle dans l'\omterm{def:s1:earnings:event}{étude d'événement} et difficile à capturer après le délai.

Une dernière variante montre la variable confondante dans toute sa gravité. On donne à chaque fonds une inclinaison propre vers un style, et les flux deviennent des paris sur les styles~: les fonds inclinés vers un style rémunéré font
bien, attirent des souscriptions nettes et achètent ce style. Le décile supérieur moins le décile inférieur des \omterm{def:s1:flows:fit}{transactions induites par les flux} gagne alors 4.9~\% sur le trimestre et 69.0~\% sur trois ans,
sans aucune inversion, car il s'agit de la prime de style, et non de la pression. Il faut contrôler un signal de flux au regard des facteurs qu'il pourrait capter par procuration.

\section{Fiches de stratégie}

\begin{strategyfile}{Inversion après vente forcée}\label{strat:s1:flows:firesale}
\sfield{Qui paie, et pourquoi} Fonds contraints de vendre par les rachats, qui acceptent un prix inférieur à la valeur pour être exécutés immédiatement.
\sfield{Instruments et plateformes} Actions détenues en commun par des fonds subissant de fortes sorties nettes.
\sfield{Signal} Ventes induites par les flux, calculées à partir des flux et des positions des fonds, dès qu'elles peuvent être connues.
\sfield{Dimensionnement et exécution} Acheter les titres les plus vendus et conserver la position tant que la pression s'inverse~; couvrir les facteurs.
\sfield{Coûts} Rotation modérée~; les titres subissent une pression vendeuse, de sorte que la liquidité est disponible pour un acheteur.
\sfield{Comment elle s'essouffle} Délais de publication qui laissent l'inversion se produire avant l'arrivée du signal~; variables confondantes avec le momentum.
\sfield{Horizon, capacité, infrastructure} De quelques mois à un an~; données de flux et de positions des fonds.
\sfield{Backtester honnêtement} Positions et flux à la date de leur publication~; contrôler les facteurs et le momentum.
\sfield{Sources} Coval et Stafford (2007)~: rendements significatifs tirés de la négociation contre des fonds contraints, 1980--2003~; simulation de ce chapitre.
\end{strategyfile}

\begin{strategyfile}{Pression induite par les flux}\label{strat:s1:flows:pressure}
\sfield{Qui paie, et pourquoi} Les investisseurs de fonds dont les flux courent après la performance passée, et les fonds tenus de les exécuter.
\sfield{Instruments et plateformes} Actions détenues par des fonds susceptibles de recevoir ou de perdre de l'argent.
\sfield{Signal} \omterm{def:s1:flows:fit}{Transactions induites par les flux} attendues~: flux prévus à partir des rendements passés des fonds, multipliés par les positions déclarées.
\sfield{Dimensionnement et exécution} Position longue sur les achats attendus, position courte sur les ventes attendues, pour le trimestre.
\sfield{Coûts} Rotation trimestrielle.
\sfield{Comment elle s'essouffle} Flux qui cessent de courir après la performance~; front-runners en encombrement (crowding) sur la même prévision.
\sfield{Horizon, capacité, infrastructure} Un trimestre~; positions, flux et rendements des fonds.
\sfield{Backtester honnêtement} Positions décalées à la date de leur publication~; modèle de flux estimé sur les seules données passées.
\sfield{Sources} Lou (2012)~: la part attendue prévoit les rendements de l'année suivante, puis s'inverse~; Coval et Stafford (2007) sur le front-running des transactions forcées.
\end{strategyfile}

\begin{strategyfile}{Flux de création-rachat d'ETF}\label{strat:s1:flows:etf}
\sfield{Qui paie, et pourquoi} Les preneurs de liquidité sur les ETF, dont les transactions atteignent le panier par arbitrage.
\sfield{Instruments et plateformes} Les ETF et leurs composantes.
\sfield{Signal} Créations et rachats quotidiens, ou prime de l'ETF par rapport à sa valeur.
\sfield{Dimensionnement et exécution} Fournir de la liquidité dans le panier face au flux~; déboucler à mesure que l'écart se résorbe.
\sfield{Coûts} Négociation de panier~; horizons courts.
\sfield{Comment elle s'essouffle} Arbitragistes plus rapides que la stratégie~; flux porteurs d'information.
\sfield{Horizon, capacité, infrastructure} Quelques jours~; fichiers de flux et de positions des ETF.
\sfield{Backtester honnêtement} Données de création avec leur heure de publication~; coûts du panier.
\sfield{Sources} Ben-David, Franzoni et Moussawi (2018)~: volatilité plus élevée et autocorrélation plus négative chez les détenteurs d'ETF.
\end{strategyfile}

\begin{strategyfile}{Signal de recouvrement des détentions}\label{strat:s1:flows:overlap}
\sfield{Qui paie, et pourquoi} Comme pour les ventes forcées~: les fonds qui détiennent les mêmes actions les vendent ensemble.
\sfield{Instruments et plateformes} Actions à détention commune concentrée.
\sfield{Signal} Le recouvrement des détenteurs d'une action avec les fonds sous tension (la mesure de recouvrement du livre~7, chapitre~28).
\sfield{Dimensionnement et exécution} En tant qu'indicateur de risque, ou comme stratégie de retournement conditionnel après une tension.
\sfield{Coûts} Faibles en tant que surcouche.
\sfield{Comment elle s'essouffle} Détentions trop anciennes pour montrer le recouvrement actuel.
\sfield{Horizon, capacité, infrastructure} Trimestres~; données de positions.
\sfield{Backtester honnêtement} Recouvrement calculé à partir des seules positions déposées.
\sfield{Sources} Coval et Stafford (2007) sur la \omterm{def:s1:flows:fit}{pression sur les prix} des titres détenus en commun~; aucun chiffre de performance vérifié.
\end{strategyfile}

\section{Tutoriel~: vendeurs forcés}\label{tut:s1:flows:tut}

\begin{tutorial}
\textbf{Objectif.} Simuler des fonds aux flux qui courent après la performance, mesurer la pression et son inversion, et négocier les flux attendus et les ventes forcées avec
le délai de publication. \textbf{État final~:} le tableau et la \cref{fig:s1:flows:path}.
\begin{tutsteps}
\item \textbf{Les fonds}~: chaque trimestre, flux issus des rendements passés, \omterm{def:s1:flows:fit}{transactions induites par les flux} à partir des positions, et \omterm{def:s1:flows:fit}{pression sur les prix} répartie sur le trimestre.
\omcode{code/firm/flowpress/firm_flowpress.py}{90}{105}{Flux, transactions induites par les flux et pression.}\label{lst:s1:flows:sim}
\item \textbf{Le livre des flux attendus}~: un modèle de flux estimé sur les trimestres passés, et des positions telles que déposées.
\omcode{code/strategies-1/09-flows/python/s1_flows.py}{100}{108}{Prévision des transactions induites par les flux du trimestre suivant.}\label{lst:s1:flows:book}
\item \textbf{Exécuter} \texttt{response}, \texttt{event\_path}, \texttt{books} (avec et sans pression) et \texttt{fig\_flows.py}.
\end{tutsteps}
\textbf{Pour aller plus loin.} Laisser la taille des fonds suivre leurs flux et observer la pression s'amplifier~; ajouter des ETF qui négocient leurs paniers mécaniquement~; raccourcir le
délai de publication et voir quelle part de l'inversion un signal plus rapide capture.
\end{tutorial}

\section{Construction~: flux des fonds et pression}\label{bld:s1:flows:flowpress}

\begin{build}
\bfield{Rôle} Positions et flux simulés des fonds, \omterm{def:s1:flows:fit}{transactions induites par les flux} et \omterm{def:s1:flows:fit}{pression sur les prix} qu'elles provoquent, pour étudier les stratégies de flux avec une vérité connue.
\bfield{Interface} \texttt{FlowConfig}, \texttt{fit(own, flow)}, \texttt{flow\_rule(trailing, beta, noise, rng)},
\texttt{simulate\_flows(R, listed, style, cfg)}.
\bfield{Règles} Flux issus de l'information disponible au début du trimestre~; pression ajoutée aux rendements sans toucher le reste du marché~; positions
enregistrées trimestre par trimestre.
\bfield{Tests d'acceptation} \texttt{code/firm/flowpress/tests/}~: \omterm{def:s1:flows:fit}{transactions induites par les flux} calculées à la main~; flux classés selon les rendements passés~; la pression suit les
transactions du trimestre, et les log-rendements ajoutés somment au niveau de pression.
\bfield{Pour aller plus loin} Tailles de fonds qui suivent les flux~; mise à l'échelle partielle des flux en transactions (fonds utilisant du cash)~; ETF.
\end{build}

\begin{omsources}
\item J.~Coval et E.~Stafford, ``Asset fire sales (and purchases) in equity markets'', \emph{Journal of Financial Economics} 86(2), 2007.
\item D.~Lou, ``A flow-based explanation for return predictability'', \emph{Review of Financial Studies} 25(12), 2012.
\item I.~Ben-David, F.~Franzoni et R.~Moussawi, ``Do ETFs increase volatility?'', \emph{Journal of Finance} 73(6), 2018.
\item X.~Gabaix et R.~S.~J.~Koijen, ``In search of the origins of financial fluctuations: the inelastic markets hypothesis'', NBER working paper 28967,
2021.
\item US SEC, Frequently Asked Questions about Form 13F.
\end{omsources}

\section{Exercices}

\begin{exercise}[$\star$]\label{exo:s1:flows:1}
Les fonds détiennent 20~\% d'une action~; ils reçoivent des entrées de 5~\% de leurs actifs. Calculer les \omterm{def:s1:flows:fit}{transactions induites par les flux} de cette action, et dire ce qu'un impact de 1.5 en fait.
\end{exercise}

\begin{exercise}[$\star$]\label{exo:s1:flows:2}
Une pression décroît avec une demi-vie de 126 jours. Calculer la part qui subsiste après 126 jours, puis après un an de 252 jours.
\end{exercise}

\begin{exercise}[$\star$]\label{exo:s1:flows:3}
Une déclaration 13F pour le trimestre clos le 31 mars doit être déposée dans les 45 jours. Calculer l'âge des positions à l'arrivée de la déclaration, puis à l'arrivée de la suivante.
\end{exercise}

\begin{exercise}[$\star\star$]\label{exo:s1:flows:4}
Pourquoi les \omterm{def:s1:flows:fit}{transactions induites par les flux} et les rendements passés désignent-ils les mêmes actions sur le marché synthétique, même sans aucune pression~?
\end{exercise}

\begin{exercise}[$\star\star$]\label{exo:s1:flows:5}
Expliquer comment des flux qui courent après la performance peuvent créer une persistance de la performance des fonds.
\end{exercise}

\begin{exercise}[$\star\star$]\label{exo:s1:flows:6}
Pourquoi le livre de l'inversion après vente forcée perd-il de l'argent, alors que l'\omterm{def:s1:earnings:event}{étude d'événement} montre une inversion~?
\end{exercise}

\begin{exercise}[$\star\star\star$]\label{exo:s1:flows:7}
\emph{Programmation.} Exécuter \texttt{event\_path(1.5, 756, 2.0)}. Expliquer pourquoi les fonds penchés transforment le signal de flux en pari de style, et comment tester un vrai signal de flux
à cet égard.
\end{exercise}

\begin{exercise}[$\star\star\star$]\label{exo:s1:flows:8}
\emph{Trouver la faille.} «~Nous calculons les \omterm{def:s1:flows:fit}{transactions induites par les flux} à partir des positions de fin de trimestre et négocions à la date de fin de trimestre~; le backtest rapporte 8~\% par an.~»
\end{exercise}

\section{Problème~: vendeurs forcés}

\begin{problem}[{Problème du week-end --- pression, retournement et délai}]\label{pb:s1:flows:1}
Les fonds simulés du chapitre sur \texttt{firm.synthmkt}, et les archives publiques.

\medskip\noindent\textbf{Partie I --- L'idée.}
\begin{enumerate}
\item Définir les \omterm{def:s1:flows:fit}{transactions induites par les flux} et la \omterm{def:s1:flows:fit}{pression sur les prix}.
\item Que trouvent Coval et Stafford~?
\item Que trouve Lou, et qu'est-ce que cela explique~?
\item Qu'est-ce que l'hypothèse des marchés inélastiques~?
\end{enumerate}

\medskip\noindent\textbf{Partie II --- La simulation.}
\begin{enumerate}[resume]
\item Décrire les fonds, leurs flux et la pression injectée.
\item Comment la pression rétroagit-elle sur les flux~?
\item Donner les réponses en coupe transversale avec et sans pression.
\item Quelle part du mouvement du trimestre s'inverse dans l'année~?
\end{enumerate}

\medskip\noindent\textbf{Partie III --- Les transactions.}
\begin{enumerate}[resume]
\item Décrire le portefeuille fondé sur les flux attendus et son résultat.
\item Décrire le portefeuille de retournement après vente forcée et expliquer pourquoi il perd.
\item Que changent les ETF~?
\item Qu'est-ce que le formulaire 13F, et quel délai de publication impose-t-il~?
\end{enumerate}

\medskip\noindent\textbf{Partie IV --- Le verdict.}
\begin{enumerate}[resume]
\item Énoncer le \emph{résultat nommé}~: la \omterm{def:s1:flows:fit}{pression sur les prix} des \omterm{def:s1:flows:fit}{transactions induites par les flux} et la part qui s'inverse dans l'année.
\item Lire l'\omterm{def:s1:earnings:event}{étude d'événement}.
\item Que montre la variante penchée vers les facteurs de style~?
\item Où obtenir des données de flux plus rapides~?
\item Comment contrôler un signal de flux vis-à-vis du momentum~?
\item Quelle fiche de stratégie faudrait-il exécuter~?
\item Qu'est-ce qui permettrait à la stratégie de vente forcée de fonctionner~?
\item En une phrase~: que vend une stratégie de flux~?
\end{enumerate}
\end{problem}

\section{Questions d'entretien}

\begin{interviewq}[$\star$ \iqroles{chercheur}]\label{iq:s1:flows:1}
Pourquoi les flux des fonds communs de placement feraient-ils bouger les cours des actions~?
\end{interviewq}

\begin{interviewq}[$\star\star$ \iqroles{chercheur, développeur}]\label{iq:s1:flows:2}
Comment calculeriez-vous les \omterm{def:s1:flows:fit}{transactions induites par les flux} pour chaque action américaine à partir de données publiques, et avec quel délai~?
\end{interviewq}

\begin{interviewq}[$\star\star$ \iqroles{chercheur}]\label{iq:s1:flows:3}
Votre signal de flux prédit les rendements. Comment montrez-vous qu'il ne s'agit pas de momentum déguisé~?
\end{interviewq}

\begin{interviewq}[$\star\star$ \iqroles{trader}]\label{iq:s1:flows:4}
Un gros fonds est connu pour être en train de liquider ses positions. Comment le négocieriez-vous, et qu'est-ce qui vous inquiéterait~?
\end{interviewq}

\begin{interviewq}[$\star\star$ \iqroles{risque}]\label{iq:s1:flows:5}
Que signifie, pour le risque d'une action, la détention commune par des fonds sous tension~?
\end{interviewq}

\begin{interviewq}[$\star\star\star$ \iqroles{chercheur}]\label{iq:s1:flows:6}
Une action est poussée de $\lambda F$ sur un trimestre, et cette poussée décroît ensuite exponentiellement avec une demi-vie de $h$ jours. Écrivez le rendement espéré d'une stratégie
qui achète à la fin du trimestre, après un délai de $L$ jours, et conserve la position pendant $H$ jours, et déterminer quand il est positif.
\end{interviewq}
